\chapter{Infinite divisibility, the eta analogue and almost generalised gamma convolutions}\label{chap:almost}\chaptermark{Infinite divisibility and almost GGCs}

Chapters~\ref{chap:basepoint} and \ref{chap:strip} attached to the reciprocal Xi function a family of laws indexed by the basepoint: the shifted clocks $T_\alpha$, obtained by subordinating the tilted centre clock, the normal scale mixtures $Y_\alpha$ and, on the Mellin side, the tilted P\'olya variables $V_\alpha$. From the basepoint one on, the shifted clocks are GGCs and the normal scale mixtures extended GGCs, with no hypothesis, and below it each of these properties is equivalent to a zero free half plane. The five sections of this chapter read these objects in other settings; the question throughout is what can be said below the basepoint one without such a half plane.

Section~\ref{ch:sato} asks what the general theory of infinitely divisible laws, in the books of Sato and of Bondesson, gives these laws, including quasi infinite divisibility of the tilted P\'olya variables above the basepoint one and its failure there. It then places the processes that the laws generate, among them a Wiener gamma integral for $T_1$ and a Riemann Ornstein Uhlenbeck process, on one map. Section~\ref{sec:wigfree} carries these laws to free probability. The Upsilon transform of Barndorff-Nielsen and Thorbj{\o}rnsen turns the centre clock into a GGC with no hypothesis, the Bercovici Pata bijection gives it a free counterpart on the half line, and RH becomes the statement that the clock is itself the Upsilon image of an infinitely divisible law. The same section reads the Jacobi parameters of the Stieltjes moments, Williams' Gamma(2) structure and the heat flow in the language of the semicircle law and of finite free convolution. Section~\ref{ch:bpeta} repeats the basepoint one calculations for the function $L(s,\chi_{12})$ of Section~\ref{ch:chi12}, whose completed function has no pole, to see which features of the example of $\zeta$ come from its pole. Section~\ref{ch:almost} gives the strongest unconditional statement that the approach delivers: every $T_\alpha$ is a GGC up to a defect carried by the zeros to the right of $\alpha$, and zero density estimates make the mass of their Cauchy kernels a vanishing fraction of the mass of the others. This falls short of a zero free half plane. Section~\ref{sec:pakes} closes the chapter with the work of Anthony Pakes on size biasing, the Lambert law and ladders of self decomposable laws: the Lambert law is a GGC with an explicit Thorin measure, by a proof that transfers to the clock except for one step, which is RH itself; the size biased clock is the clock plus an independent variable that is a mixture of exponential laws exactly under RH; and the clock lies in Urbanik's class $L_2$ of multiply self decomposable laws and not in $L_3$, with no hypothesis. None of the five sections supplies a zero free half plane.
